\chapter{Coccydynia}

\section*{Definition}
Coccydynia is pain in the coccyx (tailbone) region, typically worsened by sitting, leaning backward, or rising from a seated position. It is defined as persistent pain localized to the coccyx for more than 6 weeks without identifiable trauma in some cases.

\section*{Pathophysiology}
The coccyx consists of 3--5 fused vertebral segments connected to the sacrum via the sacrococcygeal joint. Pain arises from abnormal mobility (hypermobility or subluxation), post-traumatic periostitis, or degenerative changes of this joint. Inflammation of the surrounding ligaments and attachment of the gluteus maximus and levator ani muscles contributes to the pain. Sitting on hard surfaces increases pressure on the coccyx, aggravating symptoms.

\section*{Etiology}
\begin{itemize}
  \item \textbf{Traumatic:} Direct fall onto the buttocks, childbirth, repetitive cycling or rowing
  \item \textbf{Idiopathic:} No clear precipitating event in up to one-third of cases
  \item \textbf{Postural:} Prolonged sitting on hard surfaces, obesity (posterior pelvic tilt increases coccygeal pressure)
  \item \textbf{Biomechanical:} Sacrococcygeal hypermobility or dislocation
  \item \textbf{Neoplastic:} Chordoma, sacral tumors, metastases (rare)
  \item \textbf{Infectious:} Pilonidal cyst, sacrococcygeal abscess, osteomyelitis
\end{itemize}

\section*{Indications for Evaluation}
Seek care if:
\begin{itemize}
  \item Pain persists beyond 6--8 weeks despite conservative measures
  \item Severe pain preventing sitting or daily function
  \item Associated fever, draining sinus, or mass over the coccyx
  \item Unexplained weight loss or night pain concerning for malignancy
\end{itemize}

\section*{Related Symptoms}
\begin{itemize}
  \item Localized tenderness over the coccyx
  \item Pain exacerbated by sitting, relieved by standing or leaning forward
  \item Pain during bowel movements or sexual intercourse
  \item Referred pain to the sacrum, lower back, or buttocks
\end{itemize}
\textbf{Note:} Coccydynia is often self-limited; however, failure to improve after 8 weeks warrants imaging to exclude structural lesions.

\section*{Self-Care / Home Management}
\begin{itemize}
  \item Use of a coccyx cushion (donut or wedge pillow) when sitting
  \item Avoid prolonged sitting on hard surfaces; alternate sitting positions
  \item Ice packs to the coccyx for 15 minutes after sitting
  \item Over-the-counter NSAIDs or acetaminophen for pain
  \item Gentle stretching of the gluteal and hamstring muscles
\end{itemize}

\section*{Diagnostic Approach}
\begin{itemize}
  \item Digital rectal examination to palpate the coccyx for tenderness and mobility
  \item Dynamic radiography (sitting and standing lateral views) to assess coccygeal angle and mobility
  \item MRI if red flags present (suspected infection, mass, or after failed conservative therapy)
  \item Bone scan or CT to evaluate for occult fracture or tumor
  \item Diagnostic injection (local anesthetic with corticosteroid) to confirm the pain source
\end{itemize}

\section*{Treatment / Management Overview}
\begin{itemize}
  \item Patient education, postural modification, and cushion use as first-line
  \item NSAIDs and muscle relaxants for 2--4 weeks
  \item Physical therapy with pelvic floor relaxation and manual mobilization
  \item Local corticosteroid injection into the sacrococcygeal joint or ligament
  \item Manipulation under anesthesia for dislocated or hypermobile coccyx
  \item Coccygectomy for refractory cases after failure of all non-surgical treatments
\end{itemize}

\clearpage
